%%%%%%%%%%%%%%%%%%%%%%%%%%%%%%%%%%%%%%%%%%%%%%%%%%%%%%%%%%%%%%%%%%%%%%%%%%%%%%%%%%
\begin{frame}[fragile]\frametitle{}
\begin{center}
{\Large Examination}
\end{center}
\end{frame}

% %%%%%%%%%%%%%%%%%%%%%%%%%%%%%%%%%%%%%%%%%%%%%%%%%%%%%%%%%%%%%%%%%%%%%%%%%%%%%%%%%%
% \begin{frame}[fragile]\frametitle{}
% \begin{center}
% {\Large Param Yoga School}
% \end{center}
% \end{frame}

% %%%%%%%%%%%%%%%%%%%%%%%%%%%%%%%%%%%%%%%%%%%%%%%%%%%%%%%%%%%%%%%%%%%%%%%%%%%%%%%%%%
% \begin{frame}[fragile]\frametitle{Viva}
	% \begin{itemize}
    % \item 3 times Omkar (ॐकार), then 3 times ‘Shaanti’ (शान्ति). Sit in any sitting asana, say Sukhasana (सुखासन), eyes closed.
    % \item Prarthana (प्रार्थना): e.g. \textbf{“गुरुर्ब्रह्मा गुरुर्विष्णुः गुरुर्देवो महेश्वरः”}
	% \item Shatkarma (षट्कर्म) Etymology: "Shat" (षट्) + "Karma" (कर्म) meaning "action" or "deed."
	% \item Shatkarma in HYP (हठ योग प्रदीपिका ) in sequence: Neti (नेति ), Dhauti (धौती | वस्त्र ,  जल ,  दंड ),  Nauli (नौली), Basti (बस्ति), Kapalbhati (कपालभाति), Tratak (त्राटक)
	% \item Perform कपालभाती Kapalbhati (breathing technique)
	% \item Differences, principals, benefits of सूक्ष्म  Sukshma, स्थूल Sthula व्यायाम vyalam. Propagator?
		% \begin{itemize}
			% \item सूक्ष्म Sukshma means "subtle", slow, controlled, and mindful, one joint, Works at the micro-level, clearing subtle energy blockages (Nadis) in the body.
			% \item स्थूल Sthula means "gross", vigorous movements, multiple joints, dynamic stretches, strength-building postures, and rapid movements.
			% \item Propagator: स्वामी धीरेंद्र ब्रह्मचारी  Swami Dhirendra Brahmachari  (disciple of कार्तिकेय महाराज  Kartiyeka maharaj)
		% \end{itemize}
	% \item सूर्य नमस्कार Surya Namaskar (Sun Salutation), बीज मंत्र  beej mantra, आसन steps, Originator? पंत प्रतिनिधी औंध 
	% \item प्राणायाम  Pranayama, etymology? Prana (प्राण) meaning "life force" or "vital energy" + Ayama (आयाम) meaning "expansion" or "control."
	% \item अष्टकुम्भक  Ashtakumhak - HYP names in sequence: Surya Bhedana (सूर्यभेदन), Ujjayi (उज्जयी), Sitkari (सीतकारि ), Shitali (शीतली), Bhastrika (भस्त्रिका), Bhramari (भ्रामरी), Murcha (मूर्च्छा), Plavini (प्लाविनी)
	% \end{itemize}
% \end{frame}

% %%%%%%%%%%%%%%%%%%%%%%%%%%%%%%%%%%%%%%%%%%%%%%%%%%%%%%%%%%%%%%%%%%%%%%%%%%%%%%%%%%
% \begin{frame}[fragile]\frametitle{Viva}
% \begin{itemize}
% \item Etymology of आसन "Asana", पतंजली Patanjali's definition of आसन Asana, benefits: 
			% \begin{itemize}
			% \item "ās" (आस्), meaning "to sit" or "to be seated.
			% \item "Sthira Sukham Asanam" (स्थिरसुखमासनम्) — meaning a posture that is steady (Sthira) and comfortable (Sukham).
			% \item Benefits: Physical stability and flexibility, prepares for meditation
		% \end{itemize}			
% \item मुद्रा  Mudras (hand gestures), 5 examples: Gyan Mudra (ज्ञान मुद्रा), Chin Mudra (चिन मुद्रा), Chinmay Mudra (चिन्मय मुद्रा), Vishnu Mudra (विष्णु मुद्रा), Adi Mudra (आदि मुद्रा)
% \item बंध  Bandhas (neuro muscular locks), 4 examples, applied to redirect prana: Mula Bandha (मूल बंध), Uddiyana Bandha (उड्डियान बंध), Jalandhara Bandha (जलंधर बंध), Maha Bandha (महाबंध)
% \item Speaking on yoga topics: ज्ञान योग  jnana yoga, भक्ती योग  bhakti yoga, कर्म योग  karma yoga, अष्टांग योग  ashtang yoga, योगसूत्र yoga sutra, etc
% \end{itemize}
% \end{frame}

% %%%%%%%%%%%%%%%%%%%%%%%%%%%%%%%%%%%%%%%%%%%%%%%%%%%%%%%%%%%%%%%%%%%%%%%%%%%%%%%%%%
% \begin{frame}[fragile]\frametitle{Practical}
	% \begin{itemize}
	% \item Show proficiency in 3 Asana (आसन)
    % \item Sukshma Vyayam (सूक्ष्म व्यायाम): Greeva Shakti (ग्रीवा शक्ति) (any from I to IV) etc
    % \item Sthul Vyayam (स्थूल व्यायाम): (any Hrid (हृदय)/Sarvanga Pushti (सर्वांग पुष्टि))
	% \item Pranayam: Nadi Shodhanam (नाड़ी शोधन), Ujjayi (उज्जायी)
	% \item Perform Surya Namaskar (सूर्य नमस्कार)
	% \end{itemize}
% \end{frame}


%%%%%%%%%%%%%%%%%%%%%%%%%%%%%%%%%%%%%%%%%%%%%%%%%%%%%%%%%%%%%%%%%%%%%%%%%%%%%%%%%%
\begin{frame}[fragile]\frametitle{}
\begin{center}
{\Large Candidate reports: practical}
\end{center}
\end{frame}

%%%%%%%%%%%%%%%%%%%%%%%%%%%%%%%%%%%%%%%%%%%%%%%%%%%%%%%%%%%
\begin{frame}[fragile]\frametitle{Instructions}
    \begin{itemize}
        \item Place the mat sideways, so the camera sees the whole body in standing and lying asanas.
        \item Keep the enrolment number and an identity card ready.
        \item Check the internet connection; speak in a clear voice.
        \item If you do not know, say \textbf{``I do not know''}; do not guess. Answer only what is asked.
        \item Stay calm.
    \end{itemize}
\end{frame}

%%%%%%%%%%%%%%%%%%%%%%%%%%%%%%%%%%%%%%%%%%%%%%%%%%%%%%%%%%%
\begin{frame}[fragile]\frametitle{Preparatory Steps}
    \begin{itemize}
        \item Sit in a sitting asana such as Sukhasana (सुखासन), eyes closed. Chant Omkar (ॐकार) three times, then `Shaanti' (शान्ति) three times.
        \item Prarthana (प्रार्थना), for example: \textbf{“गुरुर्ब्रह्मा गुरुर्विष्णुः गुरुर्देवो महेश्वरः”}
        \item Sukshma Vyayam (सूक्ष्म व्यायाम): any one of Greeva Shakti (ग्रीवा शक्ति) I to IV; know who started the series and his background.
        \item Sthul Vyayam (स्थूल व्यायाम): any one of Hrid (हृदय) Gati or Sarvanga Pushti (सर्वांग पुष्टि).
    \end{itemize}
\end{frame}


%%%%%%%%%%%%%%%%%%%%%%%%%%%%%%%%%%%%%%%%%%%%%%%%%%%%%%%%%%%
\begin{frame}[fragile]\frametitle{Asanas}
    \begin{itemize}
        \item Standing asana (आसन): Trikonasana (त्रिकोणासन), hold 10 to 30 seconds, then the other side. Finish standing in Sama Sthiti (सम स्थिति).
        \item Surya Namaskar (सूर्य नमस्कार): one round, that is, the twelve positions twice, once with each leg.
        \item Lying: Shavasana (शवासन), then Viparita Karani (विपरीत करणी) (legs up, then the hips and trunk); know its cautions.
        \item Prone: Bhujangasana (भुजंगासन), then rest in Makarasana (मकरासन).
        \item Sitting: Vakrasana (वक्रासन), then the other side.
        \item A meditation asana, then Pranayama (प्राणायाम): Anuloma Viloma (अनुलोम विलोम).
        \item Bandhas (बंध): one candidate was asked Uddiyana Bandha (उड्डीयान बंध), another Jalandhara Bandha (जालंधर बंध). Uddiyana is avoided in pregnancy and menstruation.
		\item Types of meditation
			\begin{itemize}
				\item Vipassana (विपस्सना): Buddhist, taught by the Buddha
				\item Preksha (प्रेक्षा): Jain, systematised by Acharya Mahapragya from Mahavira's (महावीर) teaching
				\item Body Awareness
			\end{itemize}
    \end{itemize}
		
\end{frame}


%%%%%%%%%%%%%%%%%%%%%%%%%%%%%%%%%%%%%%%%%%%%%%%%%%%%%%%%%%%
\begin{frame}[fragile]\frametitle{Viva}
    \begin{itemize}
        \item Shatkarma (षट्कर्म), the six in order. The Hatha Pradipika order: Dhauti (धौती), Basti (बस्ति), Neti (नेति), Tratak (त्राटक), Nauli (नौली), Kapalbhati (कपालभाति). Candidates report the Gherand Samhita (घेरंड संहिता) order being asked; check it.
        \item Types of Kapalbhati, from the Gheranda Samhita: Vatakrama (वातकर्मा) (the usual form), Vyutkrama (व्यूत्कर्मा) and Sheetkrama (शीतकर्मा).
        \item Sutra Neti (सूत्र नेति): length of the thread (a candidate says 7 to 8 inches; not checked), precautions, the posture to do after it (Kagasan (कागासन)). Learned only from a qualified teacher.
        \item Asanas (आसना) asked for diabetes, as support for prevention and well-being, not cure: Vakrasana (वक्रासन), Mandukasana (मंडूकासन).
        \item Instruct Paschimottanasana (पश्चिमोत्तानासन) to the class.
        \item Give its benefits, contraindications and counter posture.
    \end{itemize}
\end{frame}

%%%%%%%%%%%%%%%%%%%%%%%%%%%%%%%%%%%%%%%%%%%%%%%%%%%%%%%%%%%%%%%%%%%%%%%%%%%%%%%%%%
\begin{frame}[fragile]\frametitle{}
\begin{center}
{\Large Candidate reports: topics}
\end{center}
\end{frame}

%%%%%%%%%%%%%%%%%%%%%%%%%%%%%%%%%%%%%%%%%%%%%%%%%%%%%%%%%%%%%%%%%%%%%%%%%%%%%%%%%%
\begin{frame}[fragile]\frametitle{Question Topics and Expectations}
    \begin{itemize}
        \item \textbf{Prayer}
        \begin{itemize}
            \item Recite a yoga prayer, such as the prayer to Patanjali or a peace chant.
            \item Know what a prayer is for, and a little of its background.
        \end{itemize}
        \item \textbf{Cleansing Techniques}
        \begin{itemize}
            \item Know Dhauti (धौती) and Neti (नेति) in full: types, method, benefits, cautions.
            \item Perform Kapalabhati (कपालभाति) and know the same details.
        \end{itemize}
        \item \textbf{Sukshma Vyayama}
        \begin{itemize}
            \item Know the Sanskrit names, for example \textit{Griva Shakti Vikasaka} (ग्रीवा शक्ति विकासक) for the neck, and every form in the syllabus.
            \item Perform every Sukshma Vyayama (सूक्ष्म व्यायाम) in the syllabus.
            \item Know the benefits and cautions.
        \end{itemize}
        \item \textbf{Sthula Vyayama}
        \begin{itemize}
            \item Know the Sanskrit names.
            \item Perform both Sthula Vyayama (स्थूल व्यायाम) practices.
            \item Know the benefits and cautions.
        \end{itemize}
        \item \textbf{Surya Namaskar}
        \begin{itemize}
            \item Know the names and the order of the positions (आसना).
            \item Know the benefits for body and mind, and the cautions.
            \item Perform it within your own capacity.
        \end{itemize}
    \end{itemize}
    {\tiny Topics reported by past candidates; not published by the Board and they may vary by agency and year. Use as practice, not as the syllabus.}
\end{frame}


%%%%%%%%%%%%%%%%%%%%%%%%%%%%%%%%%%%%%%%%%%%%%%%%%%%%%%%%%%%%%%%%%%%%%%%%%%%%%%%%%%
\begin{frame}[fragile]\frametitle{Question Topics and Expectations}
    \begin{itemize}
        \item \textbf{Asanas}
        \begin{itemize}
            \item Know every syllabus asana (आसना): technique, benefits and cautions.
            \item Perform 2 asanas the examiner names; only perform them.
            \item Perform 1 asana of your own choice and instruct it as well.
            \item Know the Sanskrit names and the reason for each name; know the counter posture.
        \end{itemize}
        \item \textbf{Breathing Practices}
        \begin{itemize}
            \item Know the 3 sectional breaths: technique, benefits, cautions.
            \item Perform and instruct them.
            \item Know the order of the full yogic breath and why it helps.
        \end{itemize}

    \end{itemize}
\end{frame}

%%%%%%%%%%%%%%%%%%%%%%%%%%%%%%%%%%%%%%%%%%%%%%%%%%%%%%%%%%%%%%%%%%%%%%%%%%%%%%%%%%
\begin{frame}[fragile]\frametitle{Question Topics and Expectations}
    \begin{itemize}
        \item \textbf{Pranayama}
        \begin{itemize}
            \item Know the meaning of \textit{Puraka} (पूरक), \textit{Rechaka} (रेचक) and \textit{Kumbhaka} (कुम्भक).
            \item Perform \textit{Anulom Vilom Pranayama} (अनुलोम विलोम प्राणायाम): hand mudra, technique, benefits, cautions.
            \item Know the beginner ratio of \textit{Puraka} to \textit{Rechaka}.
            \item \textit{Nadi Shodhana} (नाड़ी शोधन): the ratio of \textit{Puraka}, \textit{Rechaka} and \textit{Kumbhaka}; hand mudra, technique, benefits, cautions.
            \item \textit{Sheetali} (शीतली): perform it without Kumbhaka; know the Kumbhaka form, benefits and cautions.
            \item \textit{Bhramari} (भ्रामरी): perform it without Kumbhaka; know the Kumbhaka form, benefits and cautions.
            \item From the theory syllabus: how many kinds of Pranayama.
            \item Which Pranayamas the Hatha Pradipika and the Gheranda Samhita (घेरंड संहिता) name.
            \item Pranayama suited to particular conditions or age groups, for prevention and well-being.
        \end{itemize}		
    \end{itemize}
\end{frame}



%%%%%%%%%%%%%%%%%%%%%%%%%%%%%%%%%%%%%%%%%%%%%%%%%%%%%%%%%%%%%%%%%%%%%%%%%%%%%%%%%%
\begin{frame}[fragile]\frametitle{Question Topics and Expectations}
    \begin{itemize}
        \item \textbf{Bandha}
        \begin{itemize}
            \item Perform all 3 \textbf{Bandhas} (बन्ध) and instruct them.
            \item Know the technique, benefits and cautions.
            \item Theory: how many Bandhas the Hatha Yoga Pradipika (हठ योग प्रदीपिका) names.
            \item The reason for each Sanskrit name.
            \item The \textbf{Chakras} (चक्र) each Bandha is linked to in the tradition.
        \end{itemize}

        \item \textbf{Mudras}
        \begin{itemize}
            \item What a \textbf{Mudra} (मुद्रा) is, and how it is said to work.
            \item Show the syllabus Mudras and know their benefits.
            \item Mudras used in sectional breathing, and why.
        \end{itemize}

        \item \textbf{Meditation}
        \begin{itemize}
            \item Application questions: know several kinds of meditation.
            \item The eight limbs of Yoga: \textit{Dhyana} (ध्यान), \textit{Dharana} (धारणा), \textit{Samadhi} (समाधि).
            \item \textit{Pratyahara} (प्रत्याहार).
            \item The stages of \textit{Yoga Nidra} (योग निद्रा): order, benefits, cautions (care in anxiety and high blood pressure).
        \end{itemize}
    \end{itemize}
\end{frame}


%%%%%%%%%%%%%%%%%%%%%%%%%%%%%%%%%%%%%%%%%%%%%%%%%%%%%%%%%%%%%%%%%%%%%%%%%%%%%%%%%%
\begin{frame}[fragile]\frametitle{Question Topics and Expectations}
    \begin{itemize}
        \item \textbf{YOGA SESSIONS}
        \begin{itemize}
            \item Teach and speak clearly to a class.
            \item Plan a session for a person with a \textbf{disability} (विकलांग), or an \textbf{illness} (बीमार) such as \textbf{diabetes} (मधुमेह), high \textbf{blood pressure} (रक्तचाप), high \textbf{cholesterol} (कोलेस्ट्रॉल) or \textbf{back pain} (कमर दर्द); stay with prevention and well-being.
            \item Plan a session for \textbf{students} (छात्र), \textbf{women} (महिलाएँ) and other groups.
        \end{itemize}
        \item \textbf{THEORY SYLLABUS}
        \begin{itemize}
            \item Know what the theory syllabus covers.
            \item Recite the \textbf{Yoga Sutras} (योग सूत्र) 1.1 to 1.12, and \textbf{Bhagavad Gita} (भगवद गीता) verse 2.70 with its meaning; say what it means to you.
            \item Definitions of yoga in the \textbf{Yoga Vashishta} (योग वाशिष्ठ), by \textbf{Patanjali} (पतंजलि) and in the \textbf{Bhagavad Gita} (भगवद गीता).
            \item What the \textbf{Hatha Yoga Pradipika} (हठ योग प्रदीपिका) is, and its chapters; what the \textbf{Gheranda Samhita} (घेरंड संहिता) is.
            \item The idea of the \textbf{Chakras} (चक्र).
            \item \textbf{Panchkosha} (पंचकोश) and \textbf{Tridosha} (त्रिदोष).
        \end{itemize}
    \end{itemize}
\end{frame}

